% ---------------- TASK DESCRIPTION ----------------
% Wording copied exactly from the official task description
% (final_TaskDescription_ObjectToImageRatio_LearningFactory_1.pdf);
% the Student and Adviser lines are left out.
\chapter*{Task Description}
\addcontentsline{toc}{chapter}{Task Description}
\vspace*{-1.5\baselineskip}

\begin{singlespace}
\small
\setlength{\parindent}{0pt}\setlength{\parskip}{2pt}
{\centering\large\bfseries Object to image ratio and occlusion sensitivity for Material
Classification in a Learning Factory\par}
\bigskip
\noindent{\large\bfseries Motivation}\par\smallskip

Learning factory environments offer opportunities to simulate a real-world factory in a
smaller and ideal way for teaching and research purposes. The Fischertechnik learning
factory simulates a throughput of cylindrical-shaped work pieces. To that extend, an
image dataset with bounding box labels for object detection has previously been
created. The purpose of this dataset was to investigate the impact of background
elements on the object detection accuracy. That is, the dataset consists of 15
individual sets of images with the exact same parameters, only varying colour/material.
In image classification, small objects are usually a challenge as a lot of possible
complex background is present in the entire image.

\medskip
\noindent{\large\bfseries Problem Statement}\par\smallskip

This Thesis aims to systematically investigate the differences in object to image ratios
for the classification of different materials with the complex background of the
Fischertechnik learning factory. Prior to this, image classification datasets have to be
extracted from the learning factory dataset, which is labelled for object detection.
Afterwards, the research starts with implementing a CNN and occlusion sensitivity
approaches, finding baseline results and to conduct a systematic study on the object to
image ratio.
\end{singlespace}
